\section{결론}
\label{sec:conclusion}

우리는 이중 제어 설정을 도입하여 \taubench를 일반화한 \bench를 제시했으며, 조정과 의사소통 요구사항 때문에 LLM의 성능이 크게 하락함을 발견했다. 이는 사용자 요청을 해결하는 데 있어 순수 추론 능력보다 이러한 요소들이 중요한 병목임을 강조한다.

사용자 시뮬레이터를 개선하기 위해서는 아직 해야 할 일이 남아 있다. 선별된 도구로 사용자를 보강하면 치명적 오류를 피하는 데 도움이 될 수 있음을 보였지만, 이 방법을 기존 \airline 및 \retail 도메인에 어떻게 적용할 수 있을지는 아직 조사하지 않았다. 이를 수행하면 고품질 사용자 시뮬레이터를 보장하기 위한 더 일반적인 해법으로 나아갈 길이 열릴 것이다. 

벤치마크의 도메인 포괄성을 확장하는 일은 여전히 인간 전문가에 크게 의존한다. 벤치마킹 방법이 산업계에 채택되어 절실히 필요한 표준을 제공하려면, 도메인 큐레이션 과정을 자동화하는 방법을 더 깊이 조사하는 것이 중요하다.

\bench의 중요한 한계 중 하나는 대부분의 고객 지원 과제에 내재한 전문가-초보자 간극을 명시적으로 모델링하지 않는다는 점이다. 순진한 사용자와 상호작용할 때 전문가는 사용자의 멘탈 모델을 이해하고 그에 맞게 설명을 조정해야 한다. 이 간극을 메우는 AI 에이전트의 능력을 평가하고 개선하는 것은 향후 연구의 유망한 방향이며, \bench는 그러한 탐구를 위한 강력한 출발점을 제공한다.
